problem:  
Let \((M_i,\mathcal{D}_i,g_i,\mu_i)\) be a sequence of compact, equiregular sub‑Riemannian manifolds of topological dimension \(n\) and step \(s_i\ge 1\).  
Assume that each associated metric measure space \((M_i,d_{SR,i},\mu_i)\) satisfies a uniform measure contraction property \(\mathrm{MCP}(K,N)\) with fixed constants \(K\in\mathbb{R}\), \(N>1\).  
Suppose further that:  

1. The distributions \(\mathcal{D}_i\subset TM_i\) have **uniform bracket‑generation constants**, and that the corresponding nilpotent approximations at almost every point converge (in the pointed Gromov–Hausdorff sense) to the Euclidean space \((\mathbb{R}^n, d_{\mathrm{Euc}})\).  
2. The steps \(s_i\to 1\), meaning that nonholonomic directions “collapse” uniformly.  
3. The spaces \((M_i,d_{SR,i},\mu_i)\) converge in the measured Gromov–Hausdorff sense to a limit \((X,d,\mathfrak{m})\).  

**Question:**  
Do these conditions imply that the limit \((X,d,\mathfrak{m})\) satisfies the synthetic curvature‑dimension condition \(\mathrm{RCD}^*(K,N)\)?  
Equivalently, does Mosco convergence of the Cheeger energies \(\mathrm{Ch}_{M_i}\) to a quadratic limit functional on \(X\) follow from the asymptotic Riemannian behavior of the nilpotent approximations?   
If not automatically true, determine geometric–analytic necessary and sufficient conditions ensuring such **emergent infinitesimal Hilbertianity** in the limit.  

Why is it a "good" problem:  
This refined question targets the subtle boundary between the analytic and geometric aspects of synthetic curvature theory. Known results show that fixed sub‑Riemannian manifolds fail the \(\mathrm{RCD}^*\) condition due to non‑quadraticity of their Cheeger energies and non‑Hilbertian infinitesimal structure. However, what remains unexplored is the **limit process**: how and when a family of sub‑Riemannian spaces that weakly satisfy curvature conditions might converge to an \(\mathrm{RCD}^*\) space as the nonholonomic structure vanishes.  

The problem is mathematically coherent and precise enough to invite analysis through convergence of Dirichlet forms, Mosco convergence of energies, and geometric collapse of distributions. It unifies techniques from optimal transport on metric measure spaces, Riemannian approximation schemes for sub‑Riemannian structures, and analysis of nilpotent models.  

A successful resolution would yield a structural bridge between two major frameworks of modern geometry—synthetic Ricci curvature bounds and sub‑Riemannian analysis—clarifying under what degenerations Riemannian analytic behavior emerges from non‑Riemannian geometries. Simple to state but conceptually deep, it epitomizes a “narrow entrance, profound depth” problem.